% !TeX program = xelatex
% Μεταγλώττιση: xelatex lec1_exercises.tex
% Στο Overleaf: επιλέξτε XeLaTeX ως compiler.
\documentclass[11pt,a4paper]{article}

\usepackage[margin=1.9cm]{geometry}
\usepackage{fontspec}
\usepackage{polyglossia}
\setdefaultlanguage{greek}
\setotherlanguage{english}
\newfontfamily\greekfont{Libertinus Serif}
\usepackage{amsmath,amssymb,amsthm}
\usepackage{enumitem}
\usepackage{fancyvrb}
\usepackage{fancyhdr}
\usepackage[unicode,hidelinks]{hyperref}
\hypersetup{pdftitle={Το πρόβλημα του τερματισμού και τρεις μαθηματικές εικασίες}}
\newcommand{\horrule}[1]{\rule{\linewidth}{#1}}
\theoremstyle{definition}
\newtheorem{question}{Άσκηση}
\newcommand{\HALT}{\operatorname{HALT}}
\setlist[enumerate]{label=(\alph*),leftmargin=*,itemsep=3pt,topsep=5pt}
\setlength{\parindent}{0pt}
\setlength{\parskip}{5pt}
\setlength{\emergencystretch}{3em}
\pagestyle{fancy}
\fancyhf{}
\fancyfoot[L]{\small ΗΥ280 (2026) - Τσουρακάκης}
\fancyfoot[R]{\small\thepage}
\renewcommand{\headrulewidth}{0pt}
\DefineVerbatimEnvironment{pseudocode}{Verbatim}{fontsize=\small,xleftmargin=1em}

\begin{document}
\begin{center}
\horrule{0.4pt}\\[0.15cm]
{\large\textbf{ΗΥ 280 \textemdash{} Θεωρία Υπολογισμού}}\\[0.1cm]
{Ασκήσεις 1ης Διάλεξης }\\[0.1cm]
{Χαράλαμπος Ε. Τσουρακάκης}\\
\horrule{0.4pt}
\end{center}

\textbf{Κοινή υπόθεση.}
Για τις ασκήσεις 1, 2 και 3 υποθέτουμε ότι διαθέτουμε έναν ελεγκτή
$\HALT(P,D)$ που τερματίζει πάντοτε και επιστρέφει \textbf{ΝΑΙ}
αν το πρόγραμμα $P$ τερματίζει με είσοδο $D$, και \textbf{ΟΧΙ} διαφορετικά.
Γνωρίζουμε ότι τέτοιος γενικός αλγόριθμος δεν υπάρχει.
Εδώ εξετάζουμε τι θα μπορούσαμε να κάνουμε \emph{αν υπήρχε}.
Τα προγράμματά μας δεν χρειάζονται είσοδο, οπότε γράφουμε
$\HALT(P,\varepsilon)$, όπου $\varepsilon$ είναι η κενή είσοδος.
Υποθέτουμε ακριβή αριθμητική ακεραίων, χωρίς προκαθορισμένο όριο μεγέθους.

\begin{question}[Περιττοί τέλειοι αριθμοί]
Ένας θετικός ακέραιος λέγεται \textbf{τέλειος} αν ισούται με το άθροισμα
των θετικών διαιρετών του, εκτός του ίδιου. Για παράδειγμα,
\[
6=1+2+3,\qquad 28=1+2+4+7+14.
\]
Η εικασία που εξετάζουμε είναι: \textbf{δεν υπάρχει περιττός τέλειος αριθμός}.
\begin{enumerate}
\item Περιγράψτε έναν πεπερασμένο έλεγχο για το αν ένας δεδομένος θετικός
ακέραιος $n$ είναι τέλειος.
\item Γράψτε πρόγραμμα $P_1$ που τερματίζει αν και μόνο αν υπάρχει
περιττός τέλειος αριθμός. Αν βρει κάποιον, να τον εμφανίζει.
\item Εξηγήστε πώς μία κλήση στον υποθετικό ελεγκτή τερματισμού
θα προσδιόριζε αν η εικασία είναι αληθής ή ψευδής.
\end{enumerate}
\end{question}



\begin{question}[Η εικασία του Legendre]
Η εικασία του Legendre λέει ότι ανάμεσα στα τετράγωνα δύο διαδοχικών
θετικών ακεραίων υπάρχει πάντοτε τουλάχιστον ένας πρώτος αριθμός:
\[
\forall n\geq 1\;\exists p\text{ πρώτος}:\quad n^2<p<(n+1)^2.
\]
Για παράδειγμα, ανάμεσα στο $4$ και στο $9$ βρίσκεται ο πρώτος $5$,
ενώ ανάμεσα στο $9$ και στο $16$ βρίσκεται ο πρώτος $11$.
Υπενθυμίζεται ότι ένας ακέραιος $p\geq 2$ είναι πρώτος αν οι μόνοι
θετικοί διαιρέτες του είναι το $1$ και ο ίδιος.
\begin{enumerate}
\item Περιγράψτε πώς ελέγχουμε, σε πεπερασμένο χρόνο, αν η εικασία
ισχύει για έναν συγκεκριμένο θετικό ακέραιο $n$.
\item Γράψτε πρόγραμμα $P_2$ που αναζητά αντιπαράδειγμα και τερματίζει
αν και μόνο αν η εικασία είναι ψευδής.
\item Εξηγήστε πώς χρησιμοποιούμε τον υποθετικό ελεγκτή τερματισμού
για να προσδιορίσουμε αν η εικασία ισχύει.
\item Γιατί ο έλεγχος των πρώτων $10^{12}$ τιμών του $n$ δεν αποτελεί,
από μόνος του, απόδειξη της εικασίας;
\end{enumerate}
\end{question}



\begin{question}[Το πρόβλημα Brocard--Ramanujan]
Για έναν θετικό ακέραιο $n$, ορίζουμε το παραγοντικό
$n!=1\cdot2\cdots n$. Εξετάζουμε την εξίσωση
\[
n!+1=m^2,\qquad n,m\in\mathbb{Z}_{>0}.
\]
Έχουμε τις λύσεις
\[
4!+1=5^2,\qquad 5!+1=11^2,\qquad 7!+1=71^2.
\]
Η εικασία που εξετάζουμε είναι: \textbf{οι παραπάνω είναι οι μοναδικές
λύσεις σε θετικούς ακεραίους}.
\begin{enumerate}
\item Περιγράψτε πώς ελέγχουμε σε πεπερασμένο χρόνο αν ένας θετικός
ακέραιος $N$ είναι τέλειο τετράγωνο.
\item Γράψτε πρόγραμμα $P_3$ που τερματίζει αν και μόνο αν η εξίσωση
έχει λύση με $n\notin\{4,5,7\}$. Αν βρει τέτοια λύση, να την εμφανίζει.
\item Εξηγήστε τι θα σήμαινε καθεμία από τις δύο απαντήσεις του
υποθετικού ελεγκτή τερματισμού για το $P_3$.
\end{enumerate}
\end{question}



\begin{question}[Η αναπαράσταση αλλάζει το μέγεθος της εισόδου]
Στη \textbf{μοναδιαία αναπαράσταση}, ο θετικός ακέραιος $N$
γράφεται ως μια ακολουθία από $N$ μονάδες.
Για παράδειγμα, το $5$ γράφεται ως $11111$.

\begin{enumerate}
\item Γράψτε το $13$ σε μοναδιαία, δυαδική και δεκαδική
αναπαράσταση. Πόσα σύμβολα χρησιμοποιεί καθεμία;

\item Έστω $N=2^k$, όπου $k\geq 1$.
Ποιο είναι το μήκος της μοναδιαίας και ποιο της δυαδικής
αναπαράστασής του;

\item Ένας αλγόριθμος εκτελεί ακριβώς $N$ επαναλήψεις.
Εκφράστε το πλήθος των επαναλήψεων ως συνάρτηση
του μήκους της εισόδου, για καθεμία από τις δύο αναπαραστάσεις.

\item Μπορούμε να μετατρέψουμε κάθε δυαδική αναπαράσταση
μήκους $n$ στην αντίστοιχη μοναδιαία σε χρόνο πολυωνυμικό
ως προς $n$; Υποθέστε ότι σε κάθε βήμα γράφουμε το πολύ
ένα σύμβολο εξόδου.
\end{enumerate}
\end{question}



\begin{question}[Δυαδική αναπαράσταση]
Ένας θετικός ακέραιος \(N\) έχει δυαδική αναπαράσταση
με ακριβώς \(n\) bits, χωρίς αρχικά μηδενικά.

\begin{enumerate}[label=\textbf{\arabic*.}, itemsep=0.7em]
    \item Ποια είναι η μικρότερη και ποια η μεγαλύτερη
    δυνατή τιμή του \(N\);

    \item Πόσα bits απαιτούνται για την αναπαράσταση του \(2N\);
    Πώς προκύπτει από τη δυαδική αναπαράσταση του \(N\);

    \item Πότε ο αριθμός \(N+1\) χρειάζεται περισσότερα bits
    από τον \(N\);

    \item Ένας αλγόριθμος εκτελεί ακριβώς \(N\) βήματα
    με είσοδο τη δυαδική αναπαράσταση του \(N\).
    Είναι γραμμικός ως προς το \textbf{μήκος της εισόδου};
    Αιτιολογήστε την απάντησή σας.
\end{enumerate}
\end{question}

 


\begin{question}[Μία απάντηση ή μία γενική μέθοδος;]
Έστω $f:\mathbb{N}\to\{0,1\}$ μια ολική συνάρτηση.
Γράφουμε $P(x)\downarrow=y$ όταν το πρόγραμμα $P$
τερματίζει με είσοδο $x$ και επιστρέφει $y$.

Συγκρίνετε τις προτάσεις:
\[
\text{(Α)}\quad
\forall x\in\mathbb{N}\;\exists P:
P(x)\downarrow=f(x),
\]
\[
\text{(Β)}\quad
\exists P\;\forall x\in\mathbb{N}:
P(x)\downarrow=f(x).
\]

\begin{enumerate}
\item Εξηγήστε τη διαφορά τους με λόγια.
Σε ποια πρόταση μπορεί το πρόγραμμα να εξαρτάται από το $x$;

\item Δείξτε ότι η (Α) ισχύει για κάθε τέτοια συνάρτηση $f$,
ακόμη και αν η $f$ δεν είναι υπολογίσιμη.
Υπόδειξη: αρκούν δύο προγράμματα σταθερής εξόδου.

\item Ποια πρόταση εκφράζει την υπολογισιμότητα της $f$;

\item Γράψτε την άρνηση της (Β), χρησιμοποιώντας
το κατηγόρημα
\[
C(P,x):\quad
\text{«το $P(x)$ τερματίζει και επιστρέφει $f(x)$»}.
\]
Εξηγήστε τους δύο τρόπους με τους οποίους μπορεί
να αποτύχει το $C(P,x)$.
\end{enumerate}
\end{question}



\begin{question}[Τερματισμός με χρονικό όριο]
Ένας φοιτητής προτείνει τον εξής ελεγκτή:
«Προσομοιώνω το $P(D)$ για $10^{12}$ βήματα.
Αν τερματίσει, απαντώ ΝΑΙ. Διαφορετικά, απαντώ ΟΧΙ».

\begin{enumerate}
\item Γιατί αυτός ο ελεγκτής δεν λύνει το γενικό
πρόβλημα του τερματισμού;

\item Θεωρήστε τώρα το διαφορετικό πρόβλημα:
\[
\text{«Το $P(D)$ τερματίζει μέσα σε το πολύ $t$ βήματα;»}
\]
Η είσοδος περιλαμβάνει τα $P,D,t$.
Δώστε αλγόριθμο που το αποφασίζει.

\item Εξηγήστε γιατί η λύση του (β) δεν αντιφάσκει
με τη μη αποφασισιμότητα του γενικού προβλήματος.

\item Ένας δεύτερος φοιτητής προτείνει:
«Αντί για σταθερό όριο, υπολογίζω από τα $P,D$
ένα όριο $B(P,D)$.
Κάθε εκτέλεση που πρόκειται να τερματίσει
θα τερματίζει μέσα σε αυτό το όριο».

Αποδείξτε ότι δεν υπάρχει ολικός αλγόριθμος
που υπολογίζει ένα τέτοιο όριο για όλα τα $P,D$.
\end{enumerate}
\end{question}



\begin{question}[Το πρόβλημα του τερματισμού χωρίς είσοδο]
Γνωρίζουμε ότι δεν υπάρχει αλγόριθμος που να αποφασίζει,
για κάθε πρόγραμμα $P$ και είσοδο $D$,
αν το $P(D)$ τερματίζει.

Κάποιος ισχυρίζεται ότι το πρόβλημα γίνεται αποφασίσιμο
αν περιοριστούμε σε προγράμματα που δεν δέχονται είσοδο.
Υποθέτει ότι υπάρχει ελεγκτής $H_0(Q)$ που τερματίζει
πάντοτε και απαντά σωστά αν το πρόγραμμα $Q()$ τερματίζει.

\begin{enumerate}
\item Δοθέντων των $P,D$, κατασκευάστε ένα πρόγραμμα
$Q_{P,D}$ χωρίς είσοδο, το οποίο έχει ενσωματωμένα
στον κώδικά του τα $P,D$ και ικανοποιεί:
\[
Q_{P,D}()\text{ τερματίζει}
\iff
P(D)\text{ τερματίζει}.
\]

\item Εξηγήστε πώς θα χρησιμοποιούσατε το $H_0$
για να λύσετε το γενικό πρόβλημα του τερματισμού.

\item Συμπληρώστε τη σωστή κατεύθυνση της αναγωγής:
\[
\underline{\hspace{3cm}}
\;\leq\;
\underline{\hspace{3cm}}.
\]

\item Για να κατασκευάσουμε το $Q_{P,D}$,
χρειάζεται να εκτελέσουμε πρώτα το $P(D)$;
Γιατί έχει σημασία η απάντηση;
\end{enumerate}
\end{question}



\begin{question}[Μια γεύση αλγορίθμων από το Λύκειο: Ορθότητα και τερματισμός του Ευκλείδη]
Για θετικούς ακεραίους $a,b$, θεωρούμε τον αλγόριθμο:
\begin{pseudocode}
Euclid(a, b):
    όσο b ≠ 0:
        r = a mod b
        a = b
        b = r
    επέστρεψε a
\end{pseudocode}
Το $a\bmod b$ είναι το υπόλοιπο της διαίρεσης του $a$
με το $b$. Οι αναθέσεις εκτελούνται με τη σειρά που γράφονται.

\begin{enumerate}
\item Εκτελέστε τον αλγόριθμο για $(a,b)=(252,105)$,
καταγράφοντας τα διαδοχικά ζεύγη $(a,b)$.

\item Αποδείξτε ότι, για $b>0$, ο ΜΚΔ $\gcd(a,b)$  (greatest common divisor) των a,b ικανοποιεί 
\[
\gcd(a,b)=\gcd(b,a\bmod b).
\]
Υπόδειξη: γράψτε $a=qb+r$ και συγκρίνετε
τους κοινούς θετικούς διαιρέτες των δύο ζευγών.

\item Αποδείξτε ότι ο αλγόριθμος τερματίζει.

\item Συνδυάστε τα προηγούμενα για να αποδείξετε ότι
επιστρέφει τον ΜΚΔ των αρχικών αριθμών.
Γιατί το ερώτημα (β) δεν αρκεί από μόνο του;
\end{enumerate}
\end{question}



\end{document}
